\chapter{Abdominal X-Ray}

\section*{What Is This Test?}
An abdominal X-ray images your abdominal organs and structures. It's primarily used to evaluate bowel obstruction, kidney stones, foreign bodies, and perforation (free air).

\section*{Alternative Names}
AXR, KUB (kidneys, ureters, bladder), plain abdominal radiograph, abdominal radiograph


\section*{Principle}
Abdominal X-ray uses X-ray photons generated by a tungsten cathode. Photons pass through the body and are attenuated by tissues of different density. A detector (film, CR plate, or digital flat panel) captures the transmitted photons. Air appears black, fat appears dark gray, soft tissue appears light gray, bone appears white, and metal appears bright white.

\section*{History}
The abdominal X-ray was one of the first clinical applications of X-rays following Wilhelm Röntgen's discovery in 1895. It became a cornerstone of emergency radiology in the early 20th century.

\section*{Why This Test Is Ordered}
Your doctor orders an abdominal X-ray when you have severe abdominal pain, nausea and vomiting, suspected kidney stones, or when you have had recent surgery and your bowel function needs evaluation. It answers the question: is there a bowel obstruction, perforation, kidney stone, or abnormal gas pattern in your abdomen?

\section*{Is This a Primary or Secondary Test}
A primary test for initial evaluation of acute abdominal symptoms. More advanced imaging like CT is often performed afterward if the X-ray is inconclusive.

\section*{How This Test Is Performed}
You lie on a table or stand in front of the X-ray machine. Images are taken of your abdomen from the front. You may be asked to hold your breath briefly. The test takes about fifteen minutes.

\section*{What Preparation Is Needed}
Usually no preparation is needed. Tell your doctor if you are pregnant.

\section*{Contraindications and Precautions}
Pregnancy is not an absolute contraindication; the examination should be justified and optimized. Inform the team if you are pregnant or may be pregnant. There are no other specific contraindications for a plain abdominal X-ray.

\section*{Risks}
An abdominal X-ray uses a small amount of ionizing radiation. The dose is low and the benefit of diagnosing your condition typically outweighs the risk. Inform the team if you are pregnant. There are no other significant risks.

\section*{What Happens After the Test}
Results are typically available within thirty minutes. You can resume normal activities immediately.

\section*{Understanding Your Results}

\section*{Normal Results}
A normal result shows normal bowel gas pattern with no obstruction, perforation, or abnormal calcifications.
\section*{Follow-up Tests}
If the X-ray is abnormal or inconclusive, your doctor may order a CT scan of the abdomen for more detailed evaluation, a contrast-enhanced study to better visualize structures, or an ultrasound to assess specific organs.
\clearpage
